\documentclass[12pt]{article}

\usepackage[utf8]{inputenc}
\usepackage[T2A]{fontenc}
\usepackage[russian]{babel}

\usepackage{amsmath, amssymb}
\usepackage{physics}
\usepackage{geometry}
\usepackage{graphicx}
\usepackage{setspace}

\geometry{margin=2.5cm}
\onehalfspacing

\title{Продольные волны в связанной термоупругости}
\author{}
\date{}

\begin{document}

\maketitle

\section{Введение}

Волновые процессы в твёрдых телах играют ключевую роль в механике сплошных сред,
акустике, сейсмологии и неразрушающем контроле.
Особый интерес представляют продольные волны,
так как они связаны с объёмными деформациями материала
и характеризуются наибольшими скоростями распространения.

В рамках связанной термоупругости механические и тепловые поля
рассматриваются как взаимосвязанные.
Это приводит к модификации классических упругих волн
и появлению дополнительных эффектов,
связанных с тепловым расширением и теплообменом.

Цель данной статьи - описать продольные волны
в связанной термоупругой среде
и установить связь между уравнениями термоупругости
и таблицами скоростей продольных волн.

\section{Продольные упругие волны}

\subsection{Определение}

Продольной называется волна,
в которой вектор перемещений частиц среды
направлен вдоль направления распространения волны.
В результате распространения такой волны
в материале возникают чередующиеся области сжатия и разрежения.

В одномерном случае поле перемещений имеет вид
\begin{equation}
\mathbf{u} = (u(x,t), 0, 0),
\end{equation}
а деформация носит исключительно объёмный характер.

\subsection{Уравнение продольной волны в упругой среде}

Уравнение движения упругой изотропной среды записывается как
\begin{equation}
\rho \frac{\partial^2 u_i}{\partial t^2}
=
\sigma_{ij,j},
\end{equation}
где $\rho$ — плотность материала,
$\sigma_{ij}$ — тензор напряжений.

С использованием закона Гука для изотропного тела
\begin{equation}
\sigma_{ij}
=
\lambda \delta_{ij} \varepsilon_{kk}
+
2\mu \varepsilon_{ij},
\end{equation}
получается волновое уравнение для продольных колебаний:
\begin{equation}
\frac{\partial^2 u}{\partial t^2}
=
c_L^2 \frac{\partial^2 u}{\partial x^2},
\end{equation}
где
\begin{equation}
c_L = \sqrt{\frac{\lambda + 2\mu}{\rho}}
\end{equation}
— скорость продольной упругой волны.

\section{Основы связанной термоупругости}

\subsection{Связь механических и тепловых полей}

В связанной термоупругости учитывается,
что деформации вызывают изменение температуры,
а температурные изменения — дополнительные напряжения.
Основным механизмом связи является тепловое расширение.

Тензор напряжений записывается в виде
\begin{equation}
\sigma_{ij}
=
\lambda \delta_{ij} \varepsilon_{kk}
+
2\mu \varepsilon_{ij}
-
\beta \delta_{ij} (T - T_0),
\end{equation}
где $\beta$ — коэффициент термоупругой связи,
$T$ — температура,
$T_0$ — начальная температура.

\subsection{Уравнения связанной термоупругости}

Полная система уравнений включает:

\textit{Уравнение движения}:
\begin{equation}
\rho \frac{\partial^2 u_i}{\partial t^2}
=
\sigma_{ij,j}.
\end{equation}

\textit{Уравнение теплопроводности}:
\begin{equation}
\rho c \frac{\partial T}{\partial t}
=
k \nabla^2 T
+
\beta T_0 \frac{\partial \varepsilon_{kk}}{\partial t},
\end{equation}
где $c$ — удельная теплоёмкость,
$k$ — коэффициент теплопроводности.

Последний член отражает преобразование механической энергии
в тепловую.

\section{Продольные волны в связанной термоупругости}

Для продольной волны деформация определяется дивергенцией перемещений:
\begin{equation}
\varepsilon_{kk} = \frac{\partial u}{\partial x}.
\end{equation}

В результате механическое и температурное поля
становятся взаимосвязанными,
и распространение волны сопровождается
температурными возмущениями.

Связанная система приводит к существованию
модифицированной продольной волны,
скорость которой отличается от классической упругой.

\section{Скорость продольных термоупругих волн}

В предположении гармонических решений вида
\begin{equation}
u, T \sim e^{i(kx - \omega t)},
\end{equation}
получается дисперсионное соотношение,
из которого определяется скорость продольной волны:
\begin{equation}
c_L^{(te)} = \frac{\omega}{k}.
\end{equation}

В пределе слабой термоупругой связи
скорость приближённо совпадает с упругой:
\begin{equation}
c_L^{(te)} \approx \sqrt{\frac{\lambda + 2\mu}{\rho}},
\end{equation}
что объясняет использование табличных значений
скоростей продольных волн,
полученных в рамках классической упругости.

\section{Связь с таблицами скоростей продольных волн}

Таблицы скоростей продольных волн
характеризуют материал через параметры
$\lambda$, $\mu$ и $\rho$.
В термоупругой постановке эти скорости
остаются определяющими,
поскольку тепловая связь
в большинстве инженерных задач
вносит лишь малые поправки.

Таким образом,
табличные значения скоростей продольных волн
являются фундаментальной характеристикой материала
и служат основой для анализа
волновых процессов в связанной термоупругости.

\section{Заключение}

Продольные волны в связанной термоупругости
представляют собой результат взаимодействия
механических и тепловых полей.
Несмотря на наличие термоупругой связи,
основные характеристики таких волн
определяются упругими параметрами материала,
что обосновывает использование
таблиц скоростей продольных волн
при анализе термоупругих задач.

\end{document}
